\begin{apendicesenv}

\chapter{Tabela de contribuição da equipe}
\label{ap:contribuicao}

Conforme a premissa 13 do roteiro, este apêndice não conta no limite de páginas. As atividades foram levantadas dos históricos de versão das páginas da GitPage e do histórico de \textit{commits} do repositório. O percentual de esforço deve ser preenchido pela equipe a partir da planilha de esforço e totalizar 100\%.

\begin{table}[htb]
	\caption{Contribuição de cada membro na Entrega 1}
	\label{tab:contribuicao}
	\centering
	\scriptsize
	\begin{tabularx}{\textwidth}{L{3.3cm} Y C{1.6cm}}
		\toprule
		\textbf{Membro} & \textbf{Atividades realizadas} & \textbf{Esforço (\%)} \\
		\midrule
		Caio Melo Borges & Fase 1 completa: propósito e partes interessadas, descrição do software avaliado com capturas anonimizadas, modelo de qualidade adaptado, matriz Impacto~$\times$~Risco, escopo e rigor, ODS; integração da página-síntese da Fase 1; referências da Fase 1 & \pendente{--} \\
		Daniel Nunes Duarte & \pendente{Descrever as atividades (não há registro no repositório até 19/10/2026)} & \pendente{--} \\
		Davi dos Santos Brito Nobre & Criação do repositório, da estrutura MkDocs e da publicação da GitPage; primeira versão da Fase 1 (objetivo, produto, modelo e características); integração do PR da Fase 1 & \pendente{--} \\
		Felipe de Aquino Campêlo & Fase 2: visão geral e rastreabilidade, objetivos G1 a G4, questões Q1 a Q14, métricas M1 a M19, primeira versão dos níveis de pontuação e gráficos GQM & \pendente{--} \\
		Leonardo Fachinello Bonetti & Fase 2: métricas M20 e MA1, hipóteses H1 a H14 e HR1 a HR4, níveis de pontuação e critérios de julgamento, matriz de rastreabilidade, referências; integração do PR da Fase 2 & \pendente{--} \\
		\pendente{Responsável a definir} & Artigo em LaTeX (Entrega 1), questionário com estudantes e indicadores de gestão PSM/CID & \pendente{--} \\
		\midrule
		\textbf{Total} & & \textbf{100\%} \\
		\bottomrule
	\end{tabularx}
	\fonte{Elaborado pela equipe.}
\end{table}

\end{apendicesenv}
